\section{Large-content pruning on long fibers}
\label{sec:tpc29-long-content}

The $+1$ in \eqref{eq:tpc29-fiber-occupancy} is harmless when
$q=[r,s]\le J$.  In that range we can sum the $J/q$ occupancy and use
row residue energy to remove every polynomially large common content.

For $d\ge1$, define the absolute row congruence mass
\begin{equation}
 \cR_d
 :=\sum_{\alpha,\gamma}
 |\gamma_\alpha^{(1)}\gamma_\gamma^{(2)}|
 \one_{d\mid m_\alpha-m_\gamma}.
 \label{eq:tpc29-row-congruence-mass}
\end{equation}

\begin{lemma}[Row residue energy]
\label{lem:tpc29-row-residue-energy}
Uniformly for $d\ge1$,
\begin{equation}
 \boxed{
 \cR_d\ll_{\mathscr D}
 \left(\frac{Q^2}{d}+Q\right)\log(2L).}
 \label{eq:tpc29-row-residue-bound}
\end{equation}
\end{lemma}

\begin{proof}
For $i=1,2$ and $a\pmod d$, put
\[
 X_i(a)=\sum_{m_\alpha\equiv a\, (d)}|\gamma_\alpha^{(i)}|,
 \qquad
 E_i(a)=\sum_{m_\alpha\equiv a\, (d)}|\gamma_\alpha^{(i)}|^2.
\]
The distinct row integers lie in an interval of length
$O_{\mathscr D}(Q)$, so each residue class contains
$O_{\mathscr D}(Q/d+1)$ rows.  Cauchy--Schwarz gives
\[
 X_i(a)^2\ll_{\mathscr D}\left(\frac Qd+1\right)E_i(a).
\]
Hence
\begin{align*}
 \cR_d
 &=\sum_{a\, (d)}X_1(a)X_2(a)\\
 &\ll_{\mathscr D}\left(\frac Qd+1\right)
 \left(\sum_aE_1(a)\right)^{1/2}
 \left(\sum_aE_2(a)\right)^{1/2}.
\end{align*}
Apply \eqref{eq:tpc29-row-l2}.
\end{proof}

Let $\cS_{\le Y,>D}$ denote any one of the three raw ultra channels,
restricted to
\begin{equation}
 [r,s]\le Y\le J,
 \qquad
 (r,s)>D.
 \label{eq:tpc29-long-content-region}
\end{equation}
All original prefix and reflected-quotient support conditions remain
in force.

\begin{theorem}[Long-fiber large-content pruning]
\label{thm:tpc29-long-content}
Uniformly for $1\le D<Y\le J$,
\begin{equation}
 \boxed{
 |\cS_{\le Y,>D}|
 \ll_{\eps,h,\mathscr D}X^\eps
 \left\{
  \frac{XQ}{D}\log^2(2Y)+X\log^3(2Y)
 \right\}.}
 \label{eq:tpc29-long-content-bound}
\end{equation}
The same estimate holds for the sum of the three raw channels.
If $D\ge X^\eta$ for fixed $\eta>0$, then
\begin{equation}
 |\cS_{\le Y,>D}|
 \ll_{\eps,\eta,h,\mathscr D}
 XQX^{-\min\{\eta,\beta\}+\eps}.
 \label{eq:tpc29-long-content-power}
\end{equation}
\end{theorem}

\begin{proof}
Put $d=(r,s)$ and write
\[
 r=da,\qquad s=db,\qquad(a,b)=1.
\]
By target primitivity, $(d,h)=1$; by
\cref{lem:tpc29-reflected-crt}, compatibility therefore gives
$d\mid m_\alpha-m_\gamma$.  Also
\[
 [r,s]=dab.
\]
For $dab\le Y\le J$, a compatible fiber has
\begin{equation}
 O_{\mathscr D}\left(\frac{J}{dab}\right)
 \label{eq:tpc29-long-occupancy}
\end{equation}
orbit points, with the constant absorbing the endpoint $+1$.

Take absolute values and absorb the divisor coefficients and fixed
weights into $X^\eps$.  We retain the necessary condition
$d\mid m_\alpha-m_\gamma$ and the common interval support used in
\eqref{eq:tpc29-long-occupancy}, but delete the channel-specific
prefix, reflection, and squarefreeness selectors.  For each exact
$d$, the coprime parts contribute
\begin{equation}
 \sum_{\substack{a,b\ge1\\ab\le Y/d}}
 \frac1{ab}
 \ll\log^2(2Y).
 \label{eq:tpc29-harmonic-pairs}
\end{equation}
Thus \cref{lem:tpc29-row-residue-energy} gives
\begin{align}
 |\cS_{\le Y,>D}|
 &\ll X^\eps J\log^2(2Y)
 \sum_{D<d\le Y}\frac1d
 \left(\frac{Q^2}{d}+Q\right)
 \notag\\
 &\ll X^\eps\left\{
  \frac{JQ^2}{D}\log^2(2Y)
  +JQ\log^3(2Y)
 \right\}.
 \label{eq:tpc29-long-content-intermediate}
\end{align}
Use $JQ\asymp X$ to obtain
\eqref{eq:tpc29-long-content-bound}.  If $D\ge X^\eta$ and
$Q=X^{\beta+o(1)}$, the two terms save respectively $X^{-\eta}$ and
$X^{-\beta}$ relative to $XQ$; logarithms and the $o(1)$ factors are
absorbed into $X^\eps$.
\end{proof}

\begin{remark}[Two different content conditions]
\label{rem:tpc29-two-content-conditions}
For $q\le J$, the factor $J/q$ makes the simple condition
$(r,s)>X^\eta$ sufficient.  For $q\ge J$, the occupancy is only
$O(1)$ and a large gcd alone does not control the number of primitive
pairs.  The sparse theorem therefore requires the sharper condition
$q/(r,s)\le JX^{-\eta}$.
\end{remark}
